\begin{papercard}{dhondt2022}
\cf{Stages} \stagetag{H} (inference pathway that defeats a deployed fitness privacy feature), \stagetag{C} (scope and strength of a protection claim, Axis C). ACM CCS 2022 paper (KU Leuven, Ghent). Added from the gap search of 23 October 2026 as the follow-up to Hassan et al.\ \pp{hassan2018}.

\cf{Attack} After the 2018 fixes, endpoint privacy zones still leak: covered distances in activity metadata, the street grid and the points where routes enter the zone narrow the search space, and regression predicts the protected location \pc{dhondt2022}{1}{Our attack leverages distance information leaked in activity metadata, street grid data, and the locations of the entry points into the EPZ}.

\cf{Findings} On 1.4 million Strava activities the attack found the protected location for up to 85\,\% of zones \pc{dhondt2022}{1}{Our evaluation on 1.4 million Strava activities shows that our attack discovers the protected location for up to 85\% of EPZs}. Larger zones help but do not prevent it: with a 1\,km radius 55\,\% remained deanonymisable \pc{dhondt2022}{2}{even for very large radii (1 km), deanonymization remains possible for 55\% of EPZs}. Six countermeasures were proposed, each at a usability cost, and disclosed to the networks \pc{dhondt2022}{1}{We propose six countermeasures, that, however, come with a usability trade-off, and responsibly disclosed our findings and countermeasures to the major networks}; rounding reported distances was the most effective \pc{dhondt2022}{2}{We find that generalization (rounding) of reported activity distances would be the most effective countermeasure}.

\cf{Limits (card)} Attack feasibility on public activities, not a measurement of victims; no analysis of how vendors present privacy zones to users.

\cf{Use for this thesis} (1) Together with Hassan et al., a dated, quantified example that a fitness ``privacy'' feature gives partial protection: the methodology's Axis C distinguishes partial, scoped and absolute claims, and this pair shows why ``hides your home'' would be an absolute claim with a known failure condition. (2) Harm tree (fitness): data (route plus distance metadata), adversary (anyone who sees public activities), capability (location inference), consequence (physical harm, stalking). (3) The countermeasure trade-off (privacy against usability) is a fair-minded counterpoint: a vendor not adopting a countermeasure is not proof of bad faith.
\end{papercard}
